\documentclass[12pt,a4paper]{article}

% ── Packages ──────────────────────────────────────────────
\usepackage[margin=1in]{geometry}
\usepackage{graphicx}
\usepackage{amsmath}
\usepackage{booktabs}
\usepackage{array}
\usepackage{xcolor}
\usepackage{hyperref}
\usepackage{listings}
\usepackage{fancyhdr}
\usepackage{titlesec}
\usepackage{enumitem}
\usepackage{parskip}
\usepackage{multirow}
\usepackage{float}
\usepackage{caption}
\usepackage{subcaption}
\usepackage{tcolorbox}
\usepackage{fontenc}
\usepackage{inputenc}

% ── Colors ────────────────────────────────────────────────
\definecolor{cuiblue}{RGB}{0, 56, 147}
\definecolor{cuired}{RGB}{180, 20, 20}
\definecolor{codegray}{RGB}{245, 245, 245}
\definecolor{codegreen}{RGB}{0, 128, 0}
\definecolor{codepurple}{RGB}{128, 0, 128}

% ── Hyperref ──────────────────────────────────────────────
\hypersetup{
    colorlinks=true,
    linkcolor=cuiblue,
    urlcolor=cuiblue,
    citecolor=cuiblue
}

% ── Header / Footer ───────────────────────────────────────
\pagestyle{fancy}
\fancyhf{}
\rhead{\textcolor{cuiblue}{\small AI-Based Network Pathfinding \& Attack Simulation}}
\lhead{\textcolor{cuiblue}{\small CSC 262 -- Artificial Intelligence}}
\cfoot{\thepage}
\renewcommand{\headrulewidth}{0.4pt}

% ── Section Styling ───────────────────────────────────────
\titleformat{\section}{\large\bfseries\color{cuiblue}}{{\thesection.}}{0.5em}{}[\titlerule]
\titleformat{\subsection}{\normalsize\bfseries\color{cuired}}{\thesubsection.}{0.5em}{}

% ── Code Listing Style ────────────────────────────────────
\lstdefinestyle{pythonstyle}{
    backgroundcolor=\color{codegray},
    commentstyle=\color{codegreen},
    keywordstyle=\color{cuiblue}\bfseries,
    stringstyle=\color{codepurple},
    basicstyle=\ttfamily\footnotesize,
    breaklines=true,
    frame=single,
    rulecolor=\color{cuiblue},
    language=Python,
    numbers=left,
    numberstyle=\tiny\color{gray},
    showstringspaces=false
}

% ─────────────────────────────────────────────────────────
\begin{document}

% ── Title Page ────────────────────────────────────────────
\begin{titlepage}
    \centering
    \vspace*{0.5cm}

    {\Large \textbf{COMSATS University Islamabad}}\\[0.3cm]
    {\large CUI Abbottabad, Dhamtour Campus}\\[1.5cm]

    \begin{tcolorbox}[colback=cuiblue!10, colframe=cuiblue, width=\textwidth]
        \centering
        {\LARGE \textbf{AI-Based Network Pathfinding \&}}\\[0.3cm]
        {\LARGE \textbf{Attack Simulation System}}\\[0.5cm]
        {\large Project Type: Algorithm Implementation and Analysis}
    \end{tcolorbox}

    \vspace{1.2cm}

    \begin{tabular}{rl}
        \textbf{Course:}      & Artificial Intelligence (CSC 262) \\[0.4cm]
        \textbf{Instructor:}  & Zeenat Zulfiqar \\[0.4cm]
        \textbf{Email:}       & \href{mailto:zeenatzulfiqar@cuiatd.edu.pk}{zeenatzulfiqar@cuiatd.edu.pk} \\[0.4cm]
        \textbf{Section:}     & BCS-4C / BCS-8A / BCS-4B \\[0.4cm]
        \textbf{Due Date:}    & April 30, 2026 \\[0.4cm]
        \textbf{Submitted:}   & \today \\
    \end{tabular}

    \vspace{1.5cm}

    \begin{tcolorbox}[colback=gray!10, colframe=gray!50, width=0.65\textwidth]
        \centering
        \textbf{Group Members}\\[0.5cm]
        \begin{tabular}{ll}
            Name: & \underline{\hspace{5cm}} \\[0.5cm]
            Name: & \underline{\hspace{5cm}} \\[0.5cm]
            Name: & \underline{\hspace{5cm}} \\
        \end{tabular}
    \end{tcolorbox}

    \vfill
    {\small Department of Computer Science}
\end{titlepage}

% ── Table of Contents ─────────────────────────────────────
\tableofcontents
\newpage

% ═════════════════════════════════════════════════════════
\section{Introduction}
% ═════════════════════════════════════════════════════════

Modern cybersecurity landscapes are characterised by complex, interconnected network
infrastructures. Malicious threat actors do not typically compromise a target directly;
instead, they traverse a chain of intermediate systems — exploiting vulnerabilities at
each hop — until they reach a high-value asset such as a database or administrative
server. Understanding and modelling these \textit{attack paths} is fundamental to
designing resilient, secure networks.

This project bridges two disciplines: \textbf{Artificial Intelligence} and
\textbf{Cybersecurity}. A computer network is modelled as a weighted, undirected graph,
and classical as well as advanced AI search algorithms are applied to simulate an
attacker navigating from an entry point (the Internet) to a critical target (a Core
Database server). The system implements, analyses, and compares seven search
algorithms spanning uninformed, informed, local, and adversarial paradigms.

\subsection{Objectives}
\begin{itemize}
    \item Model a real-world network topology using graph data structures.
    \item Implement BFS, DFS, UCS, A*, Hill Climbing, Minimax, and Alpha-Beta Pruning from scratch.
    \item Capture and compare algorithm performance metrics: path cost, nodes expanded, and execution time.
    \item Demonstrate the failure case of Hill Climbing (local maximum entrapment).
    \item Model an adversarial attacker--defender scenario using Minimax with Alpha-Beta optimisation.
\end{itemize}

% ═════════════════════════════════════════════════════════
\section{Problem Formulation}
% ═════════════════════════════════════════════════════════

\subsection{Graph Model}

The network is represented as a weighted undirected graph $G = (V, E, W)$ where:
\begin{itemize}
    \item $V$ is the set of nodes (network entities),
    \item $E \subseteq V \times V$ is the set of edges (communication channels),
    \item $W: E \rightarrow \mathbb{R}^{+}$ assigns a vulnerability/traversal cost to each edge.
\end{itemize}

\subsection{Node Representations}

The simulated network contains \textbf{10 nodes} representing distinct network entities,
as described in Table~\ref{tab:nodes}.

\begin{table}[H]
    \centering
    \caption{Network Node Definitions}
    \label{tab:nodes}
    \begin{tabular}{llp{7cm}}
        \toprule
        \textbf{Node} & \textbf{Type} & \textbf{Description} \\
        \midrule
        Internet      & Entry Point  & Attacker's initial entry into the network \\
        Firewall      & Firewall     & Perimeter defence; filters incoming traffic \\
        WebServer     & Server       & Public-facing web application server \\
        AppServer     & Server       & Internal application logic server \\
        Workstation1  & Client       & Employee workstation (zone A) \\
        Workstation2  & Client       & Employee workstation (zone B) \\
        InternalFW    & Firewall     & Secondary internal firewall \\
        DBServer      & Database     & Intermediate database replica server \\
        AdminPC       & Client       & Administrator's privileged workstation \\
        CoreDB        & Database     & \textbf{Target:} critical core database \\
        \bottomrule
    \end{tabular}
\end{table}

\subsection{Edge Costs}

Edge weights quantify the \textit{vulnerability score} or \textit{traversal effort}
between two nodes. A lower weight indicates an easier exploitation path; a higher weight
represents a harder or more protected connection. The edge cost matrix is given in
Table~\ref{tab:edges}.

\begin{table}[H]
    \centering
    \caption{Edge Weights (Vulnerability Scores)}
    \label{tab:edges}
    \begin{tabular}{lll}
        \toprule
        \textbf{From} & \textbf{To} & \textbf{Cost} \\
        \midrule
        Internet      & Firewall      & 2 \\
        Firewall      & WebServer     & 3 \\
        Firewall      & AppServer     & 5 \\
        Firewall      & Workstation1  & 7 \\
        WebServer     & AppServer     & 2 \\
        WebServer     & InternalFW    & 6 \\
        AppServer     & InternalFW    & 3 \\
        AppServer     & Workstation1  & 4 \\
        Workstation1  & Workstation2  & 2 \\
        Workstation2  & AdminPC       & 5 \\
        Workstation2  & InternalFW    & 4 \\
        InternalFW    & DBServer      & 4 \\
        InternalFW    & AdminPC       & 3 \\
        DBServer      & CoreDB        & 2 \\
        DBServer      & AdminPC       & 3 \\
        AdminPC       & CoreDB        & 4 \\
        \bottomrule
    \end{tabular}
\end{table}

\subsection{Search Parameters}

\begin{itemize}
    \item \textbf{Start Node (Attacker Entry):} \texttt{Internet}
    \item \textbf{Goal Node (Target):} \texttt{CoreDB}
\end{itemize}

% ═════════════════════════════════════════════════════════
\section{Algorithm Implementations}
% ═════════════════════════════════════════════════════════

\subsection{Breadth-First Search (BFS)}

BFS is an uninformed search strategy that explores all nodes at the current depth level
before advancing to the next. It guarantees finding the \textbf{shortest path by number
of hops} but does not consider edge weights.

\textbf{Data Structure:} Queue (FIFO)\\
\textbf{Completeness:} Yes (finite graph)\\
\textbf{Optimality:} Yes (unweighted); No (weighted)\\
\textbf{Time Complexity:} $O(V + E)$

\begin{lstlisting}[style=pythonstyle, caption={BFS Core Logic}]
def bfs(graph, start, goal):
    queue = deque([[start]])
    visited = {start}
    while queue:
        path = queue.popleft()
        node = path[-1]
        if node == goal:
            return path, path_cost(path, graph), ...
        for neighbor, _ in graph.get(node, []):
            if neighbor not in visited:
                visited.add(neighbor)
                queue.append(path + [neighbor])
\end{lstlisting}

\subsection{Depth-First Search (DFS)}

DFS explores as deep as possible along each branch before backtracking. It uses less
memory than BFS but does not guarantee an optimal path.

\textbf{Data Structure:} Stack (LIFO)\\
\textbf{Completeness:} Yes (with visited tracking)\\
\textbf{Optimality:} No\\
\textbf{Time Complexity:} $O(V + E)$

\subsection{Uniform Cost Search (UCS)}

UCS is an informed extension of BFS that always expands the node with the
\textbf{lowest accumulated path cost}. It is guaranteed to find the optimal path in
weighted graphs.

\textbf{Data Structure:} Min-heap (priority queue)\\
\textbf{Optimality:} Yes (non-negative weights)\\
\textbf{Time Complexity:} $O((V + E) \log V)$

\subsection{A* Search Algorithm}

A* combines UCS with a heuristic function $h(n)$ to guide the search more efficiently
toward the goal. It evaluates each node using:

\begin{equation}
    f(n) = g(n) + h(n)
\end{equation}

where $g(n)$ is the accumulated cost from the start to node $n$, and $h(n)$ is an
admissible heuristic estimate of the cost from $n$ to the goal.

\subsubsection{Heuristic Design}

The heuristic used is a \textbf{scaled Euclidean distance} between the 2D canvas
positions assigned to each node:

\begin{equation}
    h(n) = \frac{\sqrt{(x_{\text{goal}} - x_n)^2 + (y_{\text{goal}} - y_n)^2}}{80}
\end{equation}

\textbf{Rationale:} Nodes are placed on a 2D coordinate plane reflecting the
network's logical topology. Nodes physically closer to CoreDB are likely to be fewer
hops away in the network. The scale factor of 80 ensures the heuristic values remain
consistent with the edge weight magnitudes, preserving admissibility (i.e., $h(n)$
never overestimates the true cost), which is required for A* to guarantee optimality.

\textbf{Optimality:} Yes (admissible heuristic)\\
\textbf{Time Complexity:} $O(E)$ in best case; exponential worst case

\subsection{Hill Climbing}

Hill Climbing is a local search algorithm. At each step, it moves to the neighbouring
node with the best (lowest) heuristic value. It does not maintain a search tree and
requires minimal memory.

\textbf{Limitation --- Local Maxima:} Hill Climbing can become trapped when no
neighbour improves the heuristic, even though the global goal is reachable via a
longer path. In our simulation, the algorithm successfully reaches the goal; however,
modifying the network to isolate certain paths demonstrates this failure mode clearly,
as indicated by the \texttt{[STUCK@node]} marker in the output.

\subsection{Minimax Algorithm}

Minimax models a two-player adversarial game:
\begin{itemize}
    \item \textbf{Attacker (Maximizer):} Selects moves to maximise progress toward \texttt{CoreDB}.
    \item \textbf{Defender (Minimizer):} Selects moves to minimise the attacker's progress.
\end{itemize}

The algorithm recursively builds a game tree to depth $d = 4$ and selects the move
with the best minimax value at each step:

\begin{equation}
    \text{minimax}(n, d, \text{isMax}) =
    \begin{cases}
        h(n) & \text{if } n = \text{goal or } d = 0 \\
        \max_{c \in \text{children}(n)} \text{minimax}(c, d-1, \text{False}) & \text{if isMax} \\
        \min_{c \in \text{children}(n)} \text{minimax}(c, d-1, \text{True}) & \text{otherwise}
    \end{cases}
\end{equation}

\subsection{Alpha-Beta Pruning}

Alpha-Beta Pruning optimises Minimax by eliminating branches of the game tree that
cannot influence the final decision:
\begin{itemize}
    \item $\alpha$: the best value the Maximizer can guarantee.
    \item $\beta$: the best value the Minimizer can guarantee.
    \item A branch is pruned when $\beta \leq \alpha$.
\end{itemize}

This reduces the effective branching factor, improving runtime without affecting the
correctness of the result. In practice, Alpha-Beta achieves the same path as Minimax
but in measurably less time.

% ═════════════════════════════════════════════════════════
\section{Results and Comparative Analysis}
% ═════════════════════════════════════════════════════════

\subsection{Empirical Results}

Table~\ref{tab:results} presents the comparative performance of all seven algorithms
on the same network scenario (Start: \texttt{Internet}, Goal: \texttt{CoreDB}).

\begin{table}[H]
    \centering
    \caption{Algorithm Comparative Analysis}
    \label{tab:results}
    \renewcommand{\arraystretch}{1.4}
    \begin{tabular}{p{2.5cm} p{6cm} r r r}
        \toprule
        \textbf{Algorithm} & \textbf{Discovered Path} & \textbf{Cost} & \textbf{Nodes} & \textbf{Time (ms)} \\
        \midrule
        BFS           & Internet$\to$Firewall$\to$WebServer$\to$InternalFW$\to$DBServer$\to$CoreDB & 17.0 & 10 & 0.047 \\
        DFS           & Internet$\to$Firewall$\to$WebServer$\to$AppServer$\to$\ldots$\to$CoreDB    & 24.0 & 10 & 0.014 \\
        UCS           & Internet$\to$Firewall$\to$AppServer$\to$InternalFW$\to$DBServer$\to$CoreDB & 16.0 & 10 & 0.070 \\
        A*            & Internet$\to$Firewall$\to$AppServer$\to$InternalFW$\to$DBServer$\to$CoreDB & 16.0 & 11 & 0.128 \\
        Hill Climbing & Internet$\to$Firewall$\to$AppServer$\to$InternalFW$\to$DBServer$\to$CoreDB & 17.0 &  5 & 0.026 \\
        Minimax       & Internet$\to$Firewall$\to$AppServer$\to$InternalFW$\to$DBServer$\to$CoreDB & 16.0 &  5 & 0.027 \\
        Alpha-Beta    & Internet$\to$Firewall$\to$AppServer$\to$InternalFW$\to$DBServer$\to$CoreDB & 16.0 &  5 & 0.021 \\
        \bottomrule
    \end{tabular}
\end{table}

\subsection{Discussion}

\subsubsection{Efficiency}
BFS and DFS expand all 10 nodes, reflecting their exhaustive nature. Hill Climbing,
Minimax, and Alpha-Beta expand only 5 nodes by following a directed greedy or
adversarial path, making them computationally lighter on this graph.

\subsubsection{Optimality}
UCS, A*, Minimax, and Alpha-Beta all discover the optimal path with cost \textbf{16}.
BFS finds a suboptimal path (cost 17) because it ignores edge weights. DFS finds the
highest-cost path (24) as it explores depth-first without cost consideration.

\subsubsection{A* vs UCS}
Both find the optimal path. A* expands slightly more nodes (11 vs 10) due to the
heuristic occasionally opening extra candidates, but on larger graphs A* would
significantly outperform UCS by directing the search toward the goal.

\subsubsection{Minimax vs Alpha-Beta}
Both produce identical results. Alpha-Beta is \textbf{$\sim$22\% faster} (0.021\,ms
vs 0.027\,ms) due to branch pruning, a difference that grows exponentially with search
depth and branching factor on more complex graphs.

\subsubsection{Hill Climbing Limitation}
Hill Climbing succeeds on this graph but is fragile: if the optimal greedy neighbour at
any step has an equal or worse heuristic than the current node, the algorithm halts and
marks the state \texttt{[STUCK@node]}. This is a known failure mode called
\textbf{local maximum entrapment}, and it cannot be resolved without backtracking
(which Hill Climbing does not support).

% ═════════════════════════════════════════════════════════
\section{User Interface}
% ═════════════════════════════════════════════════════════

A full graphical user interface (GUI) was developed using Python's \texttt{Tkinter}
library. Key features include:

\begin{itemize}
    \item \textbf{Visual Graph:} All 10 nodes rendered with colour-coded types (Entry,
    Firewall, Server, Client, Database) and edge weights labelled.
    \item \textbf{Algorithm Selector:} Dropdown menu to choose any of the 7 algorithms.
    \item \textbf{Path Highlighting:} The discovered attack path is highlighted in red
    on the graph canvas.
    \item \textbf{Metrics Table:} A live-updated results table displays path,
    cost, nodes expanded, and execution time for each run.
    \item \textbf{Run All Button:} Executes all algorithms simultaneously and populates
    the full comparison table.
\end{itemize}

The Colab notebook additionally includes a \texttt{matplotlib}-based static
visualisation for environments where Tkinter is unavailable.

% ═════════════════════════════════════════════════════════
\section{Conclusion}
% ═════════════════════════════════════════════════════════

\subsection{Summary of Findings}

This project successfully modelled a 10-node computer network as a weighted graph and
implemented seven AI search algorithms to simulate an attacker's traversal from the
Internet to a Core Database. Key findings:

\begin{itemize}
    \item \textbf{UCS and A*} produced the optimal path (cost 16) and are the most
    suitable for real-world attack path analysis where minimising exploitation effort is
    critical.
    \item \textbf{BFS} is useful for finding shortest hop-count paths but is suboptimal
    in weighted networks.
    \item \textbf{DFS} is the fastest but least reliable, often producing high-cost paths.
    \item \textbf{Hill Climbing} is efficient but susceptible to local maxima, making it
    unsuitable for production security tools.
    \item \textbf{Alpha-Beta Pruning} improves upon Minimax and is well-suited for
    modelling active attacker--defender scenarios in dynamic networks.
\end{itemize}

\subsection{Challenges}
\begin{itemize}
    \item Designing an admissible yet informative heuristic for A* required careful
    calibration of the scale factor.
    \item Implementing Minimax on a graph (rather than a traditional game tree) required
    a visited-set mechanism to prevent infinite loops in cyclic paths.
    \item Demonstrating the Hill Climbing local maximum failure required deliberate
    network design to expose the weakness.
\end{itemize}

\subsection{Future Enhancements}
\begin{itemize}
    \item Implement \textbf{Simulated Annealing} or \textbf{Genetic Algorithms} to overcome Hill Climbing limitations.
    \item Extend the adversarial model to support \textbf{dynamic edge weight changes} by the Defender in real time.
    \item Integrate with real network topology data (e.g., from NMAP scans) for practical penetration testing simulation.
    \item Add \textbf{animated path traversal} to the GUI for clearer visual demonstration.
\end{itemize}

% ═════════════════════════════════════════════════════════
\section*{References}
% ═════════════════════════════════════════════════════════
\addcontentsline{toc}{section}{References}
\begin{enumerate}
    \item Russell, S., \& Norvig, P. (2020). \textit{Artificial Intelligence: A Modern Approach} (4th ed.). Pearson.
    \item Cormen, T. H., Leiserson, C. E., Rivest, R. L., \& Stein, C. (2009). \textit{Introduction to Algorithms} (3rd ed.). MIT Press.
    \item Hart, P. E., Nilsson, N. J., \& Raphael, B. (1968). A Formal Basis for the Heuristic Determination of Minimum Cost Paths. \textit{IEEE Transactions on Systems Science and Cybernetics}, 4(2), 100--107.
    \item Stallings, W. (2017). \textit{Network Security Essentials: Applications and Standards} (6th ed.). Pearson.
\end{enumerate}

\end{document}
